% Figure: reduced pressure-temperature diagram showing the binodal
% (saturation / coexistence line) and the spinodal (locus of mechanical
% stability limit, dp/dv = 0) for the van der Waals fluid. The binodal is
% the true saturation curve located by the Maxwell equal-area construction;
% the spinodal lies inside it and marks the absolute limit of metastability.
% Between the two curves a single phase can persist metastably: a liquid
% held above its saturation pressure line is superheated, and it can survive
% up to the liquid-side spinodal before it must flash. Both curves meet at
% the critical point (Tr, pr) = (1, 1). All coordinates are dimensionless
% reduced variables computed from the reduced van der Waals equation
% pr = 8 Tr/(3 vr - 1) - 3/vr^2. The spinodal branches are clamped to
% pr >= 0 (the deep negative-pressure tension tails are not drawn).
\begin{figure}[htbp]
\centering
\begin{tikzpicture}
\begin{axis}[
  width=12.5cm, height=8.5cm,
  xlabel={reduced temperature $T_r = T/T_c$},
  ylabel={reduced pressure $p_r = p/p_c$},
  xmin=0.25, xmax=1.02, ymin=0, ymax=1.05,
  axis lines=left,
  tick label style={font=\scriptsize},
  label style={font=\small},
  legend pos=north west,
  legend style={font=\scriptsize},
  clip=false,
]
% Binodal (saturation) curve, from the Maxwell construction
\addplot[engblue, very thick] coordinates {
  (0.850,0.5045) (0.880,0.5874) (0.900,0.6470) (0.920,0.7102)
  (0.940,0.7771) (0.960,0.8476) (0.980,0.9219) (0.990,0.9605) (1.000,1.0000)
};
\addlegendentry{binodal (saturation)}
% Liquid-side spinodal branch
\addplot[engred, very thick, dashed] coordinates {
  (0.846,0.016) (0.878,0.261) (0.905,0.452) (0.927,0.601) (0.946,0.715)
  (0.961,0.802) (0.973,0.868) (0.982,0.916) (0.989,0.951) (0.994,0.975)
  (0.998,0.990) (1.000,1.000)
};
\addlegendentry{spinodal (stability limit)}
% Vapour-side spinodal branch
\addplot[engred, very thick, dashed] coordinates {
  (0.258,0.043) (0.298,0.058) (0.351,0.082) (0.427,0.125) (0.542,0.212)
  (0.625,0.293) (0.734,0.429) (0.873,0.666) (0.998,0.993)
};
% Critical point
\addplot[only marks, mark=*, engblack, mark size=2.2pt] coordinates {(1.0,1.0)};
\node[anchor=east, font=\scriptsize] at (axis cs:0.99,1.02) {critical point};
% Metastable-region label between binodal and liquid spinodal
\node[font=\scriptsize, align=center] at (axis cs:0.965,0.62)
  {superheated\\liquid\\(metastable)};
\node[font=\scriptsize, align=center] at (axis cs:0.55,0.55)
  {stable single\\phase};
\end{axis}
\end{tikzpicture}
\caption[Binodal and spinodal of the van der Waals fluid]{The binodal and
the spinodal of the van der Waals fluid in reduced coordinates, meeting at
the critical point. The binodal (solid) is the saturation curve located by
the Maxwell equal-area construction: on it, liquid and vapour coexist in
true equilibrium. The spinodal (dashed) is the locus where
$(\partial p/\partial v)_T = 0$, the absolute limit of mechanical stability.
The region between the two curves is metastable: a liquid carried above its
saturation line without nucleating a bubble is superheated and can persist
until it reaches the liquid-side spinodal, at which point no barrier remains
and it flashes to vapour. The width of that metastable gap is the room a
superheated liquid has before it must boil, and its sudden closure is the
mechanism a BLEVE exploits when a vessel wall fails and the depressurised
liquid finds itself deep inside the two-phase region.}
\label{fig:vol04ch02:spinodal-superheat}
\end{figure}
